\documentclass[11pt]{article}
\usepackage[margin=1in]{geometry}
\usepackage[T1]{fontenc}
\usepackage{lmodern,microtype,amsmath,amssymb,amsthm,mathtools,booktabs,hyperref}
\hypersetup{colorlinks=true,linkcolor=blue,urlcolor=blue,citecolor=blue,
 pdftitle={Online mistake lower bounds at arbitrary width: two nonlinear variants and a tanh obstruction},
 pdfauthor={ChatGPT-6 Astra; Kaarel Hanni; Dmitry Vaintrob}}
\newtheorem{theorem}{Theorem}[section]
\newtheorem{lemma}[theorem]{Lemma}
\newtheorem{proposition}[theorem]{Proposition}
\newtheorem{corollary}[theorem]{Corollary}
\theoremstyle{remark}\newtheorem{remark}[theorem]{Remark}
\newcommand{\sgn}{\operatorname{sign}}
\newcommand{\R}{\mathbb R}
\newcommand{\Pp}{\mathbb P}
\newcommand{\norm}[1]{\left\lVert #1\right\rVert}
\title{Online mistake lower bounds at arbitrary width:\\two nonlinear variants and a tanh obstruction}
\author{ChatGPT-6 Astra \and Kaarel H\"anni \and Dmitry Vaintrob\thanks{Authors are listed in alphabetical order.}}
\date{November 2026}
\begin{document}
\maketitle
\begin{abstract}
Can a Bayesian neural network make exponentially many online mistakes
on a simple Boolean target even when its width is arbitrarily large?
We prove that it can in two precisely specified nonlinear models.
First, the construction works with a nonlinear activation in every hidden
layer, provided that activation is sufficiently close to the identity.
The required closeness depends on the input dimension and depth but not
on width; examples include leaky ReLU with slope approaching one and a
smooth residual-tanh activation. Second, fixed standard activations---
including GELU, SiLU, ReLU, and softplus---give an exact all-width result
when only the first hidden layer is nonlinear and there are no biases.
Both models use independent Gaussian weights, a global output-margin
restriction, and exact label conditioning. A target of linear Boolean
circuit size forces exponentially many mistakes in the strong-prior
limit. At linear depth it has a teacher of width and coefficient bound
two, and the results cover every learner width above that complexity.
They concern the full posterior, including mixtures of energy minimizers.
The same near-identity learner has a polynomial iid sample bound,
independent of learner width and depth, so its online lower bound gives
a genuine separation between the two protocols.
We also construct cheap two-feature tanh competitors that invalidate a
direct lift of a scalar minimum-energy calculation. The corresponding
all-high-width question for a fixed standard activation at every hidden
layer, with a constant input coordinate, remains open.
\end{abstract}

\section{The learning question and the two results}

A bound on test error after independent training examples need not bound
the total number of errors on an unfavorable input order. Here the inputs
form a fixed sequence of distinct Boolean points. Before each label is
revealed, the learner predicts from the posterior given all earlier labels.
We ask whether a small Boolean circuit can force exponentially many errors
under this protocol, including when the learner is much wider than a
bounded-weight teacher.

The difficulty is to control the entire posterior. At small Gaussian
variance, minimum quadratic energy governs function probabilities. A
scalar network may have a unique cheapest extension of the observations,
but additional neurons can create cheaper nonlinear extensions. Thus a
scalar lower bound does not automatically survive an increase in width.

We prove two all-width variants, whose restrictions are explicit:
\begin{enumerate}
\item Every hidden layer can be nonlinear if its activation is a suitably
small perturbation of the identity. The perturbation bound is uniform in
learner width, but tends to zero as the target grows
(Theorem~\ref{thm:near}).
\item A fixed GELU, SiLU, ReLU, leaky-ReLU, or softplus activation suffices
if it occurs only in the first hidden layer, subsequent hidden layers
are linear, and the input has no appended constant coordinate
(Theorem~\ref{thm:mixed}).
\end{enumerate}
Neither theorem proves the result for one fixed standard activation in
every hidden layer of a network with a constant input coordinate.
Appendix~\ref{sec:tanh} gives a rigorous obstruction to one approach to
that remaining question; Appendix~\ref{sec:finite} gives an exact
finite-dimensional certificate of the same issue.

\subsection{Prior, conditioning, and prediction}

Depth $L$ counts all weight layers, including the linear readout, so
there are $L-1$ hidden layers. Their common learner width is $N$.
On an input cube of dimension $D$, the first matrix has $N$ rows;
its number of columns is $D+1$ when a constant coordinate is appended,
and $D$ otherwise. All Gaussian coordinates are independent, with
\begin{equation}\label{eq:prior}
 \operatorname{Var}(W^1_{ij})=\frac ND\sigma^2,\qquad
 \operatorname{Var}(W^j_{i\ell})=\sigma^2\ (2\le j<L),\qquad
 \operatorname{Var}(v_i)=\frac{\sigma^2}{N}.
\end{equation}
There are no other biases. The Gaussian density is proportional to
$\exp[-\kappa_{D,N}(\theta)/(2\sigma^2)]$, where
\begin{equation}\label{eq:energy}
 \kappa_{D,N}(\theta)=\frac DN\norm{W^1}_F^2+
       \sum_{j=2}^{L-1}\norm{W^j}_F^2+N\norm v_2^2.
\end{equation}
For $L=2$ the intermediate sum is empty.

Write $Z_\theta(\xi)$ for the real output. The prior is first conditioned
on $|Z_\theta(\xi)|>1/2$ at \emph{every} input, then on the exact observed
labels. Thus, for a history $H$ with target $f$, the posterior is the
Gaussian law conditioned on
\[
 |Z_\theta(\xi)|>1/2\quad\hbox{for all }\xi,\qquad
 f(\xi)Z_\theta(\xi)>1/2\quad\hbox{for }\xi\in H.
\]
Let $p_t$ be the posterior probability of the correct sign on round $t$.
Fresh posterior sampling makes a mistake with probability $1-p_t$;
posterior-label majority is wrong if $p_t<1/2$; the strict $2/3$
confidence criterion fails if $p_t\le2/3$.

The strong-prior limit means $\sigma\downarrow0$ at each fixed input
dimension, depth, and width. In our all-width statements, every finite
$N$ is allowed, with no upper bound, but the sufficiently-small-$\sigma$
threshold may depend on $N$. No width limit is exchanged with this limit.

For a specified architecture and depth, let $C_L(f)$ be the least
positive integer bounding both the hidden width of a strict-margin
teacher implementing $f$ and the absolute value of every teacher
coefficient. Our constructed targets have $C_L=2$ at every depth
$L\ge k+2$. Their bounded-fan-in Boolean circuit size is $\Theta(k)$.
Thus the exponential mistake count below is exponential in Boolean
circuit size, although the neural teacher depth grows with $k$.

\section{A comparator target and an exact affine minimum}
\label{sec:selector}

Fix $k\ge3$, $d=2k+1$, and $M=2^k-1$. Write an input as
$\xi=(x,y,z)$ with $x,y\in\{-1,1\}^k$ and $z\in\{-1,1\}$, and define
\[
 X=\sum_{i=0}^{k-1}2^i\frac{x_i+1}{2},\qquad
 Y=\sum_{i=0}^{k-1}2^i\frac{y_i+1}{2}.
\]
Our target is
\begin{equation}\label{eq:target}
 f(x,y,z)=+1\quad\Longleftrightarrow\quad z=-1\text{ and }X\ge Y.
\end{equation}
A bitwise comparator followed by a selector gate gives an $O(k)$-size
circuit. Every input bit is essential, so bounded fan-in also gives an
$\Omega(k)$ lower bound on circuit size.

The calibration sample consists of the following inputs, all with $z=-1$.
\begin{enumerate}
\item Positive diagonal points $x=y$: take the all-plus vector, its $k$
one-coordinate flips, and the negatives of these vectors.
\item For each $i$, take two negative carry points with
$x_i=-1,y_i=+1$ and $x_j=+1,y_j=-1$ for $j<i$. Their upper coordinates
$x_j=y_j$, $j>i$, are all plus in one point and all minus in the other.
\end{enumerate}
Discard duplicate points, leaving at most $4k+2$ examples.
For $0\le t\le M$, let $q_t$ have $(X,Y,z)=(t,0,+1)$; all these labels
are negative. Present the calibration, then
$q_k,q_{k+1},\ldots,q_M$, and finally any remaining cube points in a
fixed order. This is a permutation of the cube after duplicates in the
calibration have been removed. The designated prediction rounds will be
$q_{k+1},\ldots,q_M$, of which there are
\begin{equation}\label{eq:mistakes}
 T_k=M-k=2^k-k-1\ge2^{k-1}.
\end{equation}

\begin{lemma}[The affine optimum and a quantitative next-query gap]
\label{lem:selector}
After the calibration and $q_k,\ldots,q_t$, with $k\le t\le M$, the
unique minimum-norm affine vector with signed training margin at least
one is
\begin{equation}\label{eq:At}
 A_t(x,y,z)=-t-(t+1)z+\sum_{i=0}^{k-1}2^i(x_i-y_i).
\end{equation}
Its coefficient norm satisfies
\begin{equation}\label{eq:Qt}
 Q_t=\norm{A_t}_2^2=t^2+(t+1)^2+2\sum_{i=0}^{k-1}4^i,
 \qquad Q_M=\frac83\,4^k-2^{k+1}+\frac13.
\end{equation}
It has absolute margin at least one on the full cube. If $t<M$, then
$A_t(q_{t+1})=1$. Every affine vector $B$ satisfying the same training
constraints and $B(q_{t+1})\le-1$ has
\begin{equation}\label{eq:affine-gap}
 \norm B_2^2\ge Q_t+\frac4{d+1}.
\end{equation}
\end{lemma}
\begin{proof}
Write a competing score as $b+cz+\sum_i a_i x_i+\sum_i h_i y_i$, and put
\[
 u=b-c,\quad v=b+c,\quad w_i=(a_i-h_i)/2,\quad
 s_i=a_i+h_i,\quad U=u-1.
\]
Its squared coefficient norm is
\[
 \frac{u^2+v^2}{2}+2\sum_i w_i^2+\frac12\sum_i s_i^2.
\]
The paired diagonal examples imply $U\ge0$,
$|\sum_i s_i|\le U$, and $|\sum_j s_j-2s_i|\le U$.
Hence $|s_i|\le U$ and $\sum_i|s_i|\le kU$.
Averaging the two carry inequalities at bit $i$ gives
\[
 u-2w_i+2\sum_{j<i}w_j\le-1,
 \qquad
 w_i\ge\sum_{j<i}w_j+\frac{u+1}{2}
       \ge2^i\left(1+\frac U2\right).
\]
At $q_t$, the coordinate contribution is
$2\sum_{\mathrm{bit}_i(t)=1}w_i-\sum_{\mathrm{bit}_i(t)=0}s_i$.
The negative margin therefore implies
\[
 v\le-1-2t-tU+kU\le-(2t+1),
\]
because $t\ge k$. Combining the norm terms gives
\[
 \norm A_2^2\ge\frac{1+(2t+1)^2}{2}+2\sum_i4^i=Q_t.
\]
Equality requires $u=1$, $v=-(2t+1)$, $s_i=0$, and $w_i=2^i$,
which is exactly \eqref{eq:At}. Its scores are
$1+2(X-Y)$ for $z=-1$ and $-(2t+1)+2(X-Y)$ for $z=+1$.
They are odd integers, fit the calibration, and give
$A_t(q_j)=2j-2t-1\le-1$ for $j\le t$. This proves attainment and all
claims except \eqref{eq:affine-gap}.

The training-feasible affine vectors form a closed convex set.
Its minimum-norm point satisfies
$\langle A_t,B-A_t\rangle\ge0$ for every feasible $B$.
The extra query constraint gives
$\langle B-A_t,(1,q_{t+1})\rangle\le-2$.
Since $\norm{(1,q_{t+1})}_2=\sqrt{d+1}$, Cauchy--Schwarz implies
$\norm{B-A_t}_2^2\ge4/(d+1)$. Expanding $\norm B_2^2$ proves the gap.
\end{proof}

We will use the following elementary Gaussian principle. It also
clarifies why all our conclusions concern the full posterior.

\begin{lemma}[A strict energy gap controls posterior mass]
\label{lem:gaussian}
Fix a finite input domain and a finite-dimensional continuous network
with a positive-definite quadratic Gaussian energy $\kappa$.
Let $E$ be the event of strict global margin and a fixed training history.
Suppose a closed parameter set $B$ contains the unwanted part of $E$ and
has $\inf_B\kappa>U$, while $\inf_E\kappa\le U$.
Then $\Pp_\sigma(B\mid E)\to0$ as $\sigma\downarrow0$.
\end{lemma}
\begin{proof}
For a closed set, finite-dimensional Gaussian small-noise bounds give
\[
 \limsup_{\sigma\downarrow0}\sigma^2\log\Pp_\sigma(B)
       \le-\tfrac12\inf_B\kappa.
\]
The event $E$ is open by continuity and its strict inequalities.
Choose a point of $E$ with energy below $\inf_B\kappa$, and then a
sufficiently small ball around that point contained in $E$. The
Gaussian lower bound for this ball gives an exponentially vanishing
ratio to the conditioning-event probability. These inequalities also follow directly
by integrating the Gaussian density: a fixed small ball supplies the
lower bound, while a quadratic-energy sublevel decomposition and Gaussian
tails supply the upper bound. The energy is coercive, so a nonempty
closed constraint set attains its infimum. In particular, if all global
minimizers have one sign function, every other weak-margin sign cell has
strictly larger energy. There are finitely many such cells.
\end{proof}

Whenever $p_t\to0$ on the $T_k$ designated rounds, the claimed forms of
error follow at one common positive scale: choose $\sigma$ small enough
that every $p_t<\min\{\eta/T_k,1/3\}$. A union bound then gives
probability at least $1-\eta$ of a sampled mistake on every designated
round. Majority is wrong and the strict $2/3$ criterion fails there too.
The fixed input order does not depend on sampled predictions.

\section{A nonlinear activation at every layer, uniformly close to identity}
\label{sec:near}

This model has the appended constant coordinate $(1,\xi)$ and uses the
same activation $\phi$ in all $L-1$ hidden layers. Set $d=2k+1$,
$p=d+1$, $m=L-1$, and $Q_*=Q_M$ from \eqref{eq:Qt}.

\begin{theorem}[All widths for near-identity activations]\label{thm:near}
Fix $k\ge3$ and $L\ge2$. Choose
\begin{equation}\label{eq:epsilon}
 0<\epsilon\le
       \frac1{128m p^{3/2}Q_*^{3/2}},\qquad \alpha=1-\epsilon.
\end{equation}
Suppose $\phi(0)=0$, $\phi$ is $1$-Lipschitz, and
\begin{equation}\label{eq:close-id}
 |\phi(s)-s|\le\epsilon|s|\qquad(s\in\R).
\end{equation}
For the target and fixed input order of Section~\ref{sec:selector}, at
every fixed learner width $N\ge1$, the correct-label posterior
probability on each of the $T_k$ designated rounds tends to zero as
$\sigma\downarrow0$ under \eqref{eq:prior}. Consequently, at every
sufficiently small positive scale, posterior majority makes all $T_k$
mistakes, the strict $2/3$ confidence criterion fails on all $T_k$ rounds,
and posterior sampling makes all $T_k$ mistakes with probability at
least $1-\eta$, for any prescribed $\eta>0$.

At every depth $L\ge k+2$, the target has $C_L(f)=2$ for this activation.
Thus the lower bound applies to every $N\ge C_L(f)$ with no upper width
bound. The activation condition is independent of $N$; the positive
scale threshold may depend on $N,k,L,\phi,\eta$.
\end{theorem}

Two examples are
\[
 \phi(s)=\begin{cases}s,&s\ge0,\\(1-\epsilon)s,&s<0,\end{cases}
 \qquad\text{and}\qquad
 \phi(s)=(1-\epsilon)s+\epsilon\tanh s.
\]
Both are nonlinear for $\epsilon>0$. The second is smooth,
odd, and has derivative between $1-\epsilon$ and one. Our sufficient upper bound on $\epsilon$ in \eqref{eq:epsilon} is of
order $8^{-k}/[(L-1)k^{3/2}]$.
The theorem therefore does not give a growing-target lower bound for
one fixed standard nonlinear activation.

\begin{lemma}[Linearization on a quadratic-energy ball]
\label{lem:linearize}
Replace every hidden activation in $\theta$ by the identity, obtaining
an affine output $A_\theta\cdot(1,\xi)$. Under the theorem's hypotheses,
\begin{equation}\label{eq:linear-error}
 \max_\xi|Z_\theta(\xi)-A_\theta\cdot(1,\xi)|
 \le\epsilon m\sqrt{p/d}\left(\frac{\kappa_{d,N}(\theta)}L\right)^{L/2}.
\end{equation}
Also
\begin{equation}\label{eq:linear-energy}
 \kappa_{d,N}(\theta)\ge
       L\bigl(d\norm{A_\theta}_2^2\bigr)^{1/L}.
\end{equation}
\end{lemma}
\begin{proof}
Since $\phi$ is $1$-Lipschitz and fixes zero, it does not increase
Euclidean vector norm. Telescope through the $m$ activation locations,
replacing one at a time. The replacement error is at most $\epsilon$
times the incoming vector norm, and every remaining map has Lipschitz
constant at most its matrix Frobenius norm. Each summand is therefore
bounded by $\epsilon\sqrt p$ times the product of all $L$ matrix norms,
including the readout. The coefficients in the energy are
$d/N,1,\ldots,1,N$, whose product is $d$. Weighted
arithmetic--geometric mean gives
\[
 \prod_{j=1}^L\norm{M_j}_F
 \le d^{-1/2}\left(\frac{\kappa_{d,N}(\theta)}L\right)^{L/2},
 \qquad (M_1,\ldots,M_L)=(W^1,\ldots,W^{L-1},v).
\]
This proves \eqref{eq:linear-error}. Submultiplicativity also bounds
$\norm{A_\theta}_2$ by this product, proving \eqref{eq:linear-energy}.
\end{proof}

\begin{proof}[Proof of Theorem~\ref{thm:near}]
Fix a history through $q_t$, $k\le t<M$. First construct a cheap
consistent predictor with the wrong sign at $q_{t+1}$. Embed a scalar
path in one neuron of each hidden layer, with positive downstream
weights, whose identity-activation product vector is $P A_t$, where
\[
 P=\frac1{2\alpha^m}.
\]
Condition \eqref{eq:close-id} implies that $\phi$ preserves every
nonzero sign and $|\phi(s)|\ge\alpha|s|$. The resulting nonlinear output
has sign $A_t$ and magnitude at least
$\alpha^m P|A_t(\xi)|\ge1/2$ on the full cube. Its training signs are
correct and its next-query sign is positive.

Balance the $L$ factors against their Gaussian energy coefficients.
Explicitly, let $c_1=d/N$, $c_L=N$, and all other $c_j=1$.
Take a first row in the direction $A_t/\sqrt{Q_t}$ and choose positive
factor norms
\[
 s_j=r/\sqrt{c_j},\qquad
 r=(dP^2Q_t)^{1/(2L)}.
\]
Their product is $P\sqrt{Q_t}$, and each energy term is $r^2$.
The constructed weak-margin parameter therefore has energy
\begin{equation}\label{eq:Ut}
 U_t=L\left(\frac{dQ_t}{4\alpha^{2m}}\right)^{1/L}.
\end{equation}
Slightly increasing the readout gives an open strict-margin region of
consistent wrong predictors, with energies arbitrarily close to $U_t$.

We next show that every consistent weak-margin predictor with the
correct, negative sign at $q_{t+1}$ has energy strictly greater than
$U_t$. Suppose instead that one has energy at most $U_t$.
Lemma~\ref{lem:linearize} bounds its linearization error by
\begin{equation}\label{eq:delta}
 \delta=\frac{\epsilon m\sqrt{pQ_t}}{2\alpha^m}
             \le\frac1{128pQ_*}.
\end{equation}
Indeed, Bernoulli's inequality gives
$\alpha^m\ge1-m\epsilon\ge127/128>1/2$, and
$Q_t\le Q_*$. Put $r_0=1/2-\delta>0$.
The affine vector $B=A_\theta/r_0$ satisfies every observed unit-margin
constraint and $B(q_{t+1})\le-1$. By \eqref{eq:affine-gap} and
\eqref{eq:linear-energy},
\begin{equation}\label{eq:correct-lower}
 \kappa_{d,N}(\theta)\ge
    L\bigl(d r_0^2(Q_t+4/p)\bigr)^{1/L}.
\end{equation}
This exceeds \eqref{eq:Ut}. To check the strict inequality, set
$h=4/(pQ_t)$; here $0<h\le1/2$. Then
\[
 (1-2\delta)^2(1+h)
 \ge1+h-6\delta
 \ge1+\frac4{pQ_*}-\frac6{128pQ_*}
 >1+\frac1{16pQ_*}.
\]
On the other hand, $\epsilon\le1/2$ and $m\epsilon\le1/128$ give
\[
 \alpha^{-2m}\le\exp(4m\epsilon)\le1+8m\epsilon
             \le1+\frac1{16pQ_*}.
\]
The last inequality uses $pQ_*\ge1$ in \eqref{eq:epsilon}.
Thus $(1-2\delta)^2(Q_t+4/p)>Q_t\alpha^{-2m}$, contradicting
\eqref{eq:correct-lower} and $\kappa_{d,N}\le U_t$.

The weak-margin set fitting the history and correctly labeling the
next query is closed. At each fixed width its quadratic energy is
coercive, so its infimum is attained if the set is nonempty.
The preceding strict exclusion therefore gives an energy gap above
$U_t$. Lemma~\ref{lem:gaussian} shows that the correct-label posterior
probability tends to zero. The finite-round argument following that
lemma gives all three prediction conclusions.

For the teacher, at depth $k+2$ use $A_M/2^k$ as the first affine vector,
then $k$ scalar hidden weights equal to two, and readout one.
Every first coefficient is at most one in magnitude. Add hidden scalar
weights one if $L>k+2$. The output has the sign of $A_M$, which is $f$,
and magnitude at least $\alpha^{L-1}|A_M(\xi)|>1/2$.
Every coefficient is at most two. Zero padding is valid because
$\phi(0)=0$, hence $C_L(f)\le2$.

If $C_L(f)=1$, all layers would be scalar and all coefficients at most
one. The downstream scalar map would be $1$-Lipschitz and vanish at zero.
Its sign, up to one overall reversal, would be the first affine sign.
A strict output margin would force a first affine margin $\rho>1/2$.
Normalize that vector by $\rho$ and use the carry inequalities in
Lemma~\ref{lem:selector}. Its original top-bit antisymmetric coefficient
would exceed $2^{k-2}>1$, contradicting the coefficient bound one.
Thus $C_L(f)=2$.
\end{proof}

\begin{corollary}[A larger perturbation with exponentially many errors]
\label{cor:near-half}
Retain the architecture, target, and fixed input permutation of
Theorem~\ref{thm:near}, but count only the rounds
$q_{2^{k-1}},\ldots,q_M$. All three posterior mistake conclusions hold
on these $2^{k-1}$ rounds under the weaker sufficient condition
\begin{equation}\label{eq:epsilon-half}
 0<\epsilon\le
 \frac{2^{k+1}}{128(L-1)\sqrt{d+1}\,Q_*^{3/2}}.
\end{equation}
The teacher still has $C_L(f)=2$ for every $L\ge k+2$, and the
activation condition remains independent of learner width. The upper
bound on $\epsilon$ is of order $4^{-k}/[(L-1)\sqrt{k}]$.
\end{corollary}
\begin{proof}
Put $t_0=2^{k-1}-1$ and $a=4(t_0+1)=2^{k+1}$. For a history through
$q_t$, with $t_0\le t<M$, adding the correct negative label at
$q_{t+1}$ gives the next selector history. Its minimum affine
squared norm is exactly $Q_{t+1}$ by Lemma~\ref{lem:selector}.
Consequently the additional squared-norm cost is
\[
 \Delta_t=Q_{t+1}-Q_t=4(t+1)\ge a.
\]
This improves the coarse bound $4/p$ used in the theorem's proof.

The same scalar candidate has energy $U_t$ from \eqref{eq:Ut}.
Since $a\le Q_*$, condition \eqref{eq:epsilon-half} gives
$m\epsilon\le1/128$ and $\alpha^m>1/2$ as before. A parameter of
energy at most $U_t$ has linearization error at most
\begin{equation}\label{eq:delta-half}
 \delta=\frac{\epsilon m\sqrt{pQ_t}}{2\alpha^m}
       \le\frac{a}{128Q_*}.
\end{equation}
If it correctly labels the next query, dividing its affine
linearization by $1/2-\delta$ gives a vector of squared norm at least
$Q_{t+1}=Q_t+\Delta_t$. Set $h=\Delta_t/Q_t$.
Here $0<h<1/2$: indeed $k\ge3$, $t\ge t_0\ge3$, and
\[
 Q_t-8(t+1)\ge35+2t(t-3)>0.
\]
Thus
\[
 (1-2\delta)^2(1+h)
 \ge1+h-6\delta
 \ge1+\frac{a}{Q_*}-\frac{6a}{128Q_*}
 >1+\frac{a}{2Q_*},
\]
whereas
\[
 \alpha^{-2m}\le1+8m\epsilon
 \le1+\frac{a}{16\sqrt p\,Q_*^{3/2}}
 \le1+\frac{a}{16Q_*}.
\]
The linearization energy bound therefore excludes every correct
next-query parameter from the energy ball at $U_t$. The same closed-set
energy gap and Gaussian argument prove the posterior conclusions.
Finally $\alpha^{L-1}>1/2$ gives the identical bounded teacher, and
the exclusion of $C_L=1$ is unchanged.
\end{proof}

\subsection{Polynomial iid learning for the same nonlinear prior}
\label{sec:near-iid}

For the near-identity family, the exponential online lower bound is
accompanied by a polynomial iid upper bound for exactly the same prior,
activation, target, and width. This does not require a new restriction
on the learner's parameter support.

\begin{corollary}[An iid/online separation at every width]
\label{cor:iid}
Fix the hypotheses of Theorem~\ref{thm:near} or the larger perturbation
allowed by Corollary~\ref{cor:near-half}, including $k,L,\phi$, and
any finite width $N\ge1$. Let $\mu$ be any distribution on the input
cube. For $\varepsilon,\eta\in(0,1)$, suppose
\begin{equation}\label{eq:iid-sample}
 n\ge\frac6\varepsilon
       \left([1+d(d+1)]\log2+\log(1/\eta)\right).
\end{equation}
For every sufficiently small positive $\sigma$, with probability at least
$1-\eta$ over $n$ iid inputs labeled by $f$, the posterior gives the
correct label probability strictly greater than $2/3$ on a set of
$\mu$-measure at least $1-\varepsilon$. Its posterior-majority test
error and its mean test error under fresh posterior sampling are also
at most $\varepsilon$.

The bound on $n$ does not depend on learner width or depth. The
sufficiently-small-$\sigma$ threshold may depend on $k,L,N,\phi$ and
$\varepsilon$, but can be chosen independently of $\mu$, $n$, and the
particular target-consistent history. The activation remains a
size- and depth-dependent near-identity family.
\end{corollary}
\begin{proof}
Let $\mathcal H_d$ be the class of affine-threshold sign functions on
the cube, and set
\[
 U_*=L\left(\frac{dQ_*}{4\alpha^{2(L-1)}}\right)^{1/L}.
\]
Every weak-global-margin parameter with energy at most $U_*$ has an
affine sign function. Indeed, the linearization bound
\eqref{eq:linear-error} and the calculation \eqref{eq:delta}, now with
$t=M$, bound the error by
$1/[128(d+1)Q_*]<1/2$. Under the larger perturbation in
Corollary~\ref{cor:near-half}, the same calculation as
\eqref{eq:delta-half}, with $t=M$, instead bounds it by
$2^{k+1}/(128Q_*)\le1/128<1/2$. In either case its affine linearization
has the same sign as its nonlinear output everywhere.

The union $B$ of weak-margin cells of nonaffine Boolean functions is
closed, since the input domain is finite. It contains no point of
energy at most $U_*$. Coercivity implies $\inf_B\kappa_{d,N}>U_*$
unless $B$ is empty. On the other hand, the scalar-path construction
with $A_M$ gives the strict full-target cell
\[
 U_f=\{\theta:f(\xi)Z_\theta(\xi)>1/2\text{ for every }\xi\}
\]
energy infimum at most $U_*$. The Gaussian bounds used in
Lemma~\ref{lem:gaussian} show that
\[
 R_\sigma:=\frac{\Pp_\sigma(B)}{\Pp_\sigma(U_f)}\longrightarrow0.
\]
For every history $H$ labeled by $f$, its conditioning event $E_H$
contains $U_f$, while its nonaffine part is contained in $B$. Thus
\begin{equation}\label{eq:uniform-affine}
 \Pi_\sigma(\mathcal H_d^c\mid H)
 \le\frac{\Pp_\sigma(B)}{\Pp_\sigma(U_f)}=R_\sigma.
\end{equation}
This is one bound for all such histories, without a union bound over
histories or any dependence on the sampling distribution.

A hyperplane-arrangement count gives
\[
 |\mathcal H_d|\le2\sum_{j=0}^{d}\binom{2^d-1}{j}
                 \le2^{1+d(d+1)}.
\]
For any $\rho>0$, a fixed classifier with $\mu$-error greater than
$\rho$ remains consistent with $n$ iid labels with probability at most
$e^{-n\rho}$. A union bound implies that, with probability at least
$1-\eta$, every consistent affine classifier has test error at most
$\rho$, provided
\[
 n\ge\frac{[1+d(d+1)]\log2+\log(1/\eta)}\rho.
\]
Choose $\rho=\varepsilon/6$ and choose $\sigma$ small enough that
$R_\sigma\le\varepsilon/6$. On this sample event, the mean posterior
error at a fresh input is at most
$\rho+R_\sigma\le\varepsilon/3$.
Write $q(\xi)$ for the posterior probability of the wrong label.
Posterior-majority errors are contained in $\{q\ge1/2\}$, and failures
of strict $2/3$ confidence are exactly $\{q\ge1/3\}$.
Markov's inequality bounds their $\mu$-measures by
$2\varepsilon/3$ and $\varepsilon$, respectively, proving the result.
\end{proof}

Thus the same nonlinear learner has an
$O((d^2+\log(1/\eta))/\varepsilon)$ iid sample guarantee and at least
$2^k-k-1$ designated online errors, or $2^{k-1}$ under the larger
perturbation allowance, with $d=2k+1$, at every fixed width.
For fixed $k,L,N$, one sufficiently small positive scale satisfies both
claims. The contrast is between different ways of choosing examples;
it does not require changing the prior during learning.

\section{Fixed standard activations with one nonlinear layer}
\label{sec:mixed-activation}

An exact extension is available for several standard nonlinear
activations if only the first hidden layer is nonlinear and the input
has no constant coordinate. These are explicit changes to the learner:
the result below does not concern a network with a nonlinear activation
at every hidden layer, and does not allow free first-layer biases.

Let $x\in\{-1,1\}^p$, let $N\ge2$, and write
\[
 Z_\theta(x)=vW^{L-1}\cdots W^2\phi(W^1x),\qquad L\ge2.
\]
The entries of $W^1$, $W^2,\ldots,W^{L-1}$, and $v$ have independent
centered Gaussian laws with variances $(N/p)\sigma^2$, $\sigma^2$,
and $\sigma^2/N$, respectively. Their quadratic energy is
\[
 \kappa_N(\theta)=\frac pN\|W^1\|_F^2+
 \sum_{j=2}^{L-1}\|W^j\|_F^2+N\|v\|_2^2.
\]
Condition on $|Z_\theta(x)|>1/2$ at every Boolean input,
and then on the observed labels.

Assume that $\phi$ is continuous and
\begin{equation}\label{eq:mixed-odd}
 \phi(t)-\phi(-t)=ct\quad(t\in\mathbb R),\qquad c>0,
\end{equation}
and assume in addition that either $t\phi(t)>0$ for $t\ne0$, or
$\phi$ is nondecreasing. GELU, SiLU, ReLU, and softplus satisfy these
conditions with $c=1$: here GELU is $t\Phi(t)$ with $\Phi$ the
standard normal distribution function, SiLU is $t/(1+e^{-t})$, ReLU is
$\max\{t,0\}$, and softplus is $\log(1+e^t)$. The first two preserve
the sign of nonzero inputs; the latter two are nondecreasing. A leaky ReLU with negative slope
$\alpha\in(0,1]$ satisfies them with $c=1+\alpha$.

\begin{lemma}[Exact energy for a paired history]
\label{lem:mixed-energy}
Let $H\subset\{-1,1\}^p\times\{-1,1\}$ be a nonempty history closed
under replacing $(x,y)$ by $(-x,-y)$. Assume the homogeneous linear margin constraints
are feasible. Let $A_*$ be the unique minimum-Euclidean-norm vector
satisfying $yA_*\cdot x\ge1/2$ on $H$. Suppose also that
$|A_*\cdot x|\ge1/2$ on the full cube. Then the minimum energy among
parameters with weak full-cube margin $1/2$ that fit $H$ is
\[
 E_*=L\left(\frac{4p\|A_*\|_2^2}{c^2}\right)^{1/L}.
\]
Every minimizing parameter represents $\operatorname{sign}(A_*\cdot x)$.
The actual conditioned Gaussian posterior therefore concentrates on
this sign function as $\sigma\downarrow0$, at each fixed $N,L$.
\end{lemma}
\begin{proof}
Put $b=vW^{L-1}\cdots W^2$. Equation~\eqref{eq:mixed-odd} gives
\[
 \frac{Z_\theta(x)-Z_\theta(-x)}2=A\cdot x,
 \qquad A=\frac c2 bW^1.
\]
The paired training constraints imply $yA\cdot x\ge1/2$, hence
$\|A\|_2\ge\|A_*\|_2$. Frobenius submultiplicativity and
arithmetic--geometric mean give
\begin{align*}
 \kappa_N(\theta)
 &\ge L\left(p\|W^1\|_F^2
       \prod_{j=2}^{L-1}\|W^j\|_F^2\|v\|_2^2\right)^{1/L}\\
 &\ge L\left(p\|W^1\|_F^2\|b\|_2^2\right)^{1/L}
 \ge L\left(4p\|A\|_2^2/c^2\right)^{1/L}\ge E_*.
\end{align*}

For attainment, let $u=A_*/\|A_*\|_2$ and put
\[
 e_1=(1,-1,0,\ldots,0)/\sqrt2.
\]
Choose unit vectors $e_2,\ldots,e_{L-1}\in\mathbb R^N$ and set
\[
 r=\left(4p\|A_*\|_2^2/c^2\right)^{1/(2L)},\qquad
 s_1=r\sqrt{N/p},\quad s_j=r\ (2\le j<L),\quad
 s_L=r/\sqrt N.
\]
Take $W^1=s_1e_1u^\top$, $W^j=s_je_je_{j-1}^\top$, and
$v=s_Le_{L-1}^\top$, with the evident interpretation when $L=2$.
The opposite first-layer pair cancels the even part of $\phi$, so
\[
 Z_\theta(x)=\frac c2\left(\prod_{j=1}^Ls_j\right)u\cdot x
             =A_*\cdot x.
\]
Each energy summand equals $r^2$, proving attainment.

At any energy minimizer all the preceding inequalities are equalities.
Thus $A=A_*$, and equality in
$\|bW^1\|_2\le\|b\|_2\|W^1\|_F$ forces
$W^1=\lambda b^\top u^\top$ for some $\lambda>0$. Consequently
\[
 Z_\theta(x)=F(u\cdot x),\qquad
 F(s)=\sum_i b_i\phi(\lambda b_is).
\]
If $t\phi(t)>0$ for $t\ne0$, each nonzero summand has the sign of
$s$, which proves the sign claim. If $\phi$ is nondecreasing, then
$F$ is nondecreasing, including when some $b_i$ are negative.
At least one training constraint for $A_*$ is active, since otherwise
$A_*$ could be scaled down. Its antipodal pair yields
\[
 F\left(\frac1{2\|A_*\|_2}\right)\ge\frac12,\qquad
 F\left(-\frac1{2\|A_*\|_2}\right)\le-\frac12.
\]
Every cube projection has absolute value at least
$1/(2\|A_*\|_2)$, so monotonicity again proves the sign claim.

There are finitely many Boolean sign functions, their weak-margin
parameter cells are closed, and the energy is coercive. All minimizers
have the same sign function, so every other consistent sign cell has a
strictly larger energy infimum. Increasing the readout in the attained
linear construction by a factor tending to one makes all margins
strict without changing the limiting energy. Lemma~\ref{lem:gaussian} now proves the posterior claim, including
any mixture of tied minimizers.
\end{proof}

\begin{theorem}[All widths for one nonlinear layer]\label{thm:mixed}
For every $k\ge3$, there is a Boolean target on $p=2k+2$ inputs,
computable by an $O(k)$-size bounded-fan-in Boolean circuit, and a fixed
sequence of distinct inputs forcing at least
$2^k-k-1\ge2^{k-1}$ mistakes in the strong-prior limit of the preceding
learner, at every fixed $N\ge2$ and $L\ge2$.
For each of the listed activations, the target's bounded-width-and-weight
complexity is exactly two at every depth $L\ge k+2$.
Thus the lower bound holds at every width at least that complexity.
For every $\eta>0$, all sufficiently small positive $\sigma$ make
posterior sampling wrong on every designated round with probability at
least $1-\eta$; posterior majority is wrong and the strict $2/3$
confidence criterion fails on all those rounds. The threshold may depend
on $N,k,L,\phi,\eta$; the input order does not depend on $N$ or $\sigma$.
\end{theorem}
\begin{proof}
Write $M=2^k-1$, add a variable Boolean coordinate $s$ to the selector
construction of Section~\ref{sec:selector}, and use the homogeneous target
\[
 F(s,x,y,z)=\operatorname{sign}\left(
   -Ms-(M+1)z+\sum_{i=0}^{k-1}2^i(x_i-y_i)\right).
\]
Equivalently, $F(s,x,y,z)=s f(sx,sy,sz)$ for the target in
\eqref{eq:target}. Coordinatewise multiplication by $s$ and the final
sign reversal cost only $O(k)$ Boolean gates.
The restriction $s=1$ is the comparator $f$, so the circuit size
is in fact $\Theta(k)$. Lift every history
point from Section~\ref{sec:selector} to $s=1$ and immediately follow it by its antipode, with the
opposite label. Lemma~\ref{lem:selector} gives, after the
paired history through query $t$, the unique unit-margin optimum
\[
 \widehat A_t=-ts-(t+1)z+\sum_i2^i(x_i-y_i).
\]
Its scores are odd nonzero integers on the full cube. Apply
Lemma~\ref{lem:mixed-energy} with $A_*=\widehat A_t/2$. The first point of each
next query pair, for $t=k,\ldots,M-1$, receives the incorrect positive
sign with posterior probability tending to one. The antipodal points
need not be counted. They are distinct and preserve a paired history
before each counted query. A finite-round union bound gives the
positive-scale sampled-prediction statement.

At depth $L=k+2$, take first-layer rows $\pm\widehat A_M/2^k$. In the first
linear hidden layer use the first row $(2/c,-2/c)$ and a zero second
row. Each of the remaining $k-1$ linear hidden layers multiplies the
first coordinate by two and keeps the second zero. Use readout
$(1,0)$. The resulting output is exactly $\widehat A_M$, with absolute margin at
least one. All coefficients are at most two for the listed activations,
and the width is two. Additional identity layers allow every
$L\ge k+2$.

Complexity one would give scalar layers with all coefficients bounded
by one. The odd projection, normalized from margin $1/2$ to unit
margin, would then have every coefficient bounded by $c\le2$.
Restricting its paired unit-margin constraints to $s=1$, the carry
inequalities in Lemma~\ref{lem:selector} instead force its highest-bit
coefficient difference divided by two to be at least
$2^{k-1}>2$. This contradiction proves that the complexity is exactly
two.
\end{proof}

Free first-layer biases remain a substantive obstacle. For example,
when $\beta\ge0$,
\[
 \operatorname{ReLU}(\beta+t)-\operatorname{ReLU}(\beta-t)
 =t+\operatorname{clip}(t,-\beta,\beta),
\]
so the odd projection is nonlinear. Additional nonlinear hidden layers
also destroy the linear odd-projection identity used above.


\section{What remains open}

The two theorems remove the upper bound on learner width, but do so in
different restricted models. The first retains a nonlinear activation
at every hidden layer and a constant input coordinate; its activation
becomes increasingly close to linear as the target grows. The second
uses a fixed standard activation, but removes biases and makes every
subsequent hidden layer linear. In both cases an explicit identity or
energy gap controls all low-energy competitors, which is what permits
a conclusion about the actual posterior rather than a selected minimizer.

For fixed tanh, ReLU, GELU, or a fixed-slope leaky ReLU at every hidden
layer, with an appended constant input, we have not proved or disproved
an exponential online mistake lower bound at every sufficiently large
width. In particular, the tanh constructions in the appendices do not
prove that a wider learner makes few mistakes. They prove that certain
scalar implementations cease to be energy minimizers after widening;
other wider implementations of wrong functions might still be cheaper.

These questions concern cumulative errors on a prescribed order.
The near-identity family has the polynomial iid upper bound in
Corollary~\ref{cor:iid}; no iid sample-complexity lower bound is asserted.

\appendix
\section{A two-feature tanh obstruction to lifting the scalar argument}
\label{sec:tanh}

We now return to tanh in \emph{every} hidden layer, an appended constant
input, and the prior \eqref{eq:prior}. The following is an energy
comparison, not a statement about the predictions of the wider posterior.
It shows why the scalar affine minimum cannot simply be inserted as
the minimum energy of the wider architecture.

\begin{proposition}[A cheap classifier consistent with the entire suffix]
\label{prop:tanh-shortcut}
Let $k\ge128$, $d=2k+1$, and $L=d+1=2k+2$.
For every width $N\ge2$, the Boolean function
\[
 g(x,y,z)=+1\quad\Longleftrightarrow\quad z=-1\text{ and }x_0=y_0
\]
has a strict-global-margin tanh implementation with every coefficient
bounded by two and
\[
 \kappa_{d,N}\le L+89+12\log d<2L.
\]
It fits every calibration example and every negative query $q_t$ from
Section~\ref{sec:selector}.
In contrast, every scalar tanh implementation of $\sgn A_t$, for
$k\le t\le M$, with weak global output margin $1/2$ has energy at least
$2L$.
\end{proposition}
\begin{proof}
Every positive calibration point is diagonal. At every negative carry
point, $x_0\ne y_0$: either bit zero is the highest differing bit, or
it is one of the lower carry bits. Every $q_t$ has $z=+1$.
Thus $g$ fits the entire stated sequence.

Put
\[
 \lambda=(82d)^{-1/2},\qquad
 s=x_0-y_0+\frac72+\frac52z,\qquad
 r=\frac{\tanh(4\lambda)}{\tanh(2\lambda)}
   =\frac2{1+\tanh^2(2\lambda)}.
\]
Consider the two-feature combination
\[
 F(s)=\tanh(2\lambda s)-r\tanh(\lambda s).
\]
Writing $a=\tanh(\lambda s)$ and $b=\tanh(2\lambda)$, the double-angle
identity gives
\begin{equation}\label{eq:tanh-two}
 F(s)=\frac{2a(b^2-a^2)}{(1+a^2)(1+b^2)}.
\end{equation}
It is odd, positive for $0<s<2$, and negative for $s>2$.
The possible $s$ values are $-1,1,3$ for $z=-1$ and $4,6,8$ for
$z=+1$. Exactly $s=1$ has the positive target label, so $\sgn F=g$.

For $N\ge2$, put $n=\lfloor N/2\rfloor$. Use $n$ copies of each
first-layer feature and zero in any remaining neuron. In each of the
$N$ second-layer neurons, assign coefficients $1/n$ to the first
feature group and $-r/n$ to the second. Every second-layer activation
is then $\tanh(F(s))$.
The squared affine coefficient norm of $s$ is $41/2$.
The first-layer energy is at most $82d\lambda^2=1$.
Since $N/n\le3$ and $r\le2$, the second-layer energy is at most
$(N/n)(1+r^2)\le15$.
A later layer with all entries $a/N$ maps a common activation $h$ to
$\tanh(ah)$ and has Frobenius energy $a^2$. A readout with all entries
$1/N$ has weighted energy one and returns that common activation.

We need a margin estimate to choose the remaining layers.
For $d\ge7$, $8\lambda<1/2$. On $[0,1/2]$ we have
$\tanh'(u)>1/2$ and $\tanh u\ge u/2$.
For $s=1$, these imply $a\ge\lambda/2$,
$b-a\ge\lambda/2$, and $b+a\ge3\lambda/2$.
Equation~\eqref{eq:tanh-two}, with denominator at most four, gives
$F(1)\ge3\lambda^3/16$. For $s\in\{3,4,6,8\}$, use
$a\ge3\lambda/2$, $a-b\ge\lambda/2$, and $a+b\ge5\lambda/2$
to get $|F(s)|\ge15\lambda^3/16$. Oddness handles $s=-1$.
Every second-layer activation therefore has absolute value at least
\[
 h_0=\tanh(3\lambda^3/16)\ge3\lambda^3/32.
\]

Any $J\le d$ subsequent gain-one tanh layers retain at least $h_0/2$.
Indeed, $\tanh u\ge u-u^3/3$ for $u\ge0$. Iterating from $h_0$,
all virtual values are at most $h_0$, so the total loss is at most
$Jh_0^3/3<h_0/2$, using $h_0\le3\lambda^3/16$ and $J\le d$. Monotonicity transfers this bound to larger actual
activation magnitudes. For $0<h\le1/2$, concavity and
$\tanh1>3/4$ give $\tanh(2h)>(3/2)h$.
Once $h\ge1/2$, one more gain-two layer makes it strictly greater than
$1/2$ and subsequent such layers preserve that property. Hence
\[
 m=\left\lceil\log_{3/2}\!\left[\frac{32}3(82d)^{3/2}\right]
       \right\rceil+1
\]
final gain-two layers suffice after any preceding gain-one layers.
For $d\ge257$, elementary logarithm estimates give
$m\le25+4\log d<d-2=L-3$, so the available depth suffices.
Use gain one in the other $L-3-m$ layers. The total energy is at most
\[
 1+15+(L-3-m)+4m+1
 =L+14+3m\le L+89+12\log d<2L.
\]
All physical weights are bounded by two.

For the scalar lower bound, absorb the scalar downstream signs into the
first affine vector, using oddness of tanh. Let $\rho>0$ be the least
absolute first preactivation. Dividing its affine vector by $\rho$
gives a unit-margin representation of $\sgn A_t$. Its squared norm is
at least $Q_t$, since this sign function fits the history in
Lemma~\ref{lem:selector}. Since $|\tanh u|\le|u|$, the product of
$\rho$ and all remaining scalar weight magnitudes must be at least
$1/2$. Weighted arithmetic--geometric mean yields
\[
 \kappa_{d,1}\ge L(dQ_t/4)^{1/L}.
\]
But $Q_t\ge(2/3)(4^k-1)$ and
$d(1-4^{-k})>24$ for the present $k$, so
$dQ_t/4\ge4^{k+1}=2^L$. The lower bound is at least $2L$.
\end{proof}

The target comparator itself still has a bounded-weight tanh teacher:
start with $A_M/2^k$, apply tanh, then use $2k$ further scalar hidden
weights two and readout one. The smallest first activation is at least
$(3/4)2^{-k}$. Until it exceeds $1/2$, each following layer multiplies
it by more than $3/2$; after $2k$ layers the resulting lower estimate is
$(3/4)(9/8)^k>1/2$. All coefficients are at most two. The carry
inequalities exclude complexity one exactly as in the near-identity
teacher argument. Thus $C_L(f)=2$ at this depth.

Equal replication of a scalar tanh chain to any width preserves its
energy: duplicate the first row, divide every subsequent scalar
connection and readout by the width, and use the coefficients in
\eqref{eq:energy}. Proposition~\ref{prop:tanh-shortcut} therefore shows
that these replicated scalar implementations are not minimizers in the
larger architecture. It does \emph{not} lower-bound every wider
implementation of $\sgn A_t$, and does not identify the wider posterior's
prediction. Some different wider implementation of a wrong function
could still have lower energy than the displayed competitor.

% Standalone insertion; requires amsmath, amsthm and a proposition environment.
\section{An exact finite-width tanh certificate}
\label{sec:finite}

The following calculation identifies a concrete obstruction to extending a
width-one online-mistake proof merely by replicating its neurons. It does
not establish the prediction of the unrestricted wider-network posterior:
a different wider implementation of a wrong function could still have
smaller energy.

Use $d=7$ input bits $(x_0,x_1,x_2,y_0,y_1,y_2,z)$, seven tanh hidden
layers of common width $N$, and a linear readout. There is one appended
constant input and no other biases. The Gaussian quadratic energy is
\[
 \kappa_N=\frac7N\|W^1\|_F^2+
       \sum_{\ell=2}^{7}\|W^\ell\|_F^2+N\|v\|_2^2.
\]
Put $X=\sum_{i=0}^2 2^i(x_i+1)/2$, $Y=\sum_{i=0}^2 2^i(y_i+1)/2$ and
\[
 f(x,y,z)=+1\ \Longleftrightarrow\ z=-1\text{ and }X\ge Y.
\]
After the calibration in Section~\ref{sec:selector} and queries
$X=3,4,5,6$, $Y=0$, $z=+1$, the affine minimum is the sign function
\[
 g_6=\operatorname{sign}\!\left(-6-7z+
                       \sum_{i=0}^2 2^i(x_i-y_i)\right).
\]
Its minimum squared affine coefficient norm at unit global margin is
$Q(g_6)=42+36+49=127$. In particular, it wrongly labels the next query
$(X,Y,z)=(7,0,+1)$ positive.

\begin{proposition}
For every integer $N\ge64$, the target $f$ has a width-$N$ implementation
with all weights of absolute value less than three, global output margin
greater than $1/2$, and
\[
 \kappa_N<15.7.
\]
Every width-one chain computing $g_6$ with global output margin at least
$1/2$ has energy at least
\[
 8(7\cdot127/4)^{1/8}>15.7.
\]
Equal replication of that scalar chain to width $N$ preserves its energy.
Thus no equally replicated scalar minimizer for this history minimizes
energy among all width-$N$ networks.
\end{proposition}
\begin{proof}
Let $n=\lfloor N/64+1/2\rfloor$ and $p=n/N$. For $N\ge64$,
\[
 1/95\le p\le1/48.
\]
Indeed, when $n=1$ we have $64\le N\le95$; when $n\ge2$, the inequalities
$64n-32\le N\le64n+31$ imply the displayed bounds.

Write
\[
 U=1+\sum_{i=0}^2 2^i(x_i-y_i),\quad
 \alpha=\frac{55}{100},\quad b=\frac{17}{100},\quad
 r=\frac{147}{1000},\quad s=\frac{675}{1000}.
\]
The first $n$ neurons of layer one compute $a=\tanh(\alpha U)$;
the remaining $N-n$ compute $c=\tanh(b(z+1))$. Every neuron of layer two
computes the same value $h_2=\tanh(ra-sc)$, using weights $r/n$ and
$-s/(N-n)$ on the two groups. The next five layers are equal replicas
with respective effective scalar weights
\[
 (w_3,w_4,w_5,w_6,w_7)
   =\frac1{1000}(1361,1365,1373,1389,1419).
\]
Each physical weight in layer $j$ is $w_j/N$. The readout has $N$ equal
weights $(3/2)/N$.

When $z=-1$, $c=0$ and $|U|\ge1$, so the second-layer preactivation
has the sign of $U$ and absolute value at least $r\tanh\alpha$.
When $z=+1$, it is at most $r-s\tanh(2b)<0$.
Consequently the network has the target's sign everywhere, and the
minimum second-layer preactivation magnitude is at least
\[
 m=\min\{r P(\alpha),\ sP(2b)-r\},\qquad P(u)=u-u^3/3.
\]
Here $\tanh u\ge P(u)$ for $u\ge0$: the difference vanishes at zero
and has derivative $u^2-\tanh^2u\ge0$.
Starting with $\underline h_2=P(m)$ and recursively setting
$\underline h_j=P(w_j\underline h_{j-1})$ gives positive rational
lower bounds on the hidden magnitudes. Exact rational arithmetic gives
\[
 \frac32\underline h_7>\frac12
 \qquad\text{(numerically, }0.5033684798\ldots\text{)}.
\]
The accompanying script \texttt{verification/check\_finite\_certificate.py}
verifies this inequality without floating-point arithmetic.

Since the squared affine coefficient norm of $U$ is $43$, the energy is
exactly
\[
 E(p)=7\bigl(43p\alpha^2+2(1-p)b^2\bigr)
       +\frac{r^2}{p}+\frac{s^2}{1-p}
       +\sum_{j=3}^7w_j^2+\frac94.
\]
This is convex in $p$. Exact rational evaluation at the two endpoints
shows
\[
 E(1/95)<15.666,\qquad E(1/48)<15.590,
\]
so $E(p)<15.7$ on the entire interval. The largest physical weight in
this implementation is $4\alpha=2.2$.

For the scalar lower bound, let $\rho$ be its least absolute first-layer
preactivation. Its first coefficient vector has squared norm at least
$127\rho^2$. Since $|\tanh u|\le|u|$, the product of $\rho$ and the
absolute values of its remaining seven weights is at least $1/2$.
Applying arithmetic--geometric mean to the eight energy terms yields
$\kappa_1\ge8(7\cdot127/4)^{1/8}$. The strict comparison with $15.7$
is the rational inequality $(157/80)^8<889/4$.
\end{proof}


\paragraph{Verification status.}
The main arguments and the tanh constructions have conventional proofs
and independent AI-agent mathematical checks. The finite certificate in
Appendix~\ref{sec:finite} uses exact rational arithmetic. These results
are not Lean formalized. The accompanying verification folder identifies
the checked source snapshot and the scope of the reviews.
\end{document}
